\chapter{תוצרים וקריטריוני קבלה}\label{ch:deliv}

כפי שהוסבר בתחילת הדוח (״היקף: חזון מול מטלה״, פרק~\ref{ch:exec}), התוצרים בטבלה~\ref{tab:deliv} הם של החברה המתוכננת. מה שנבנה בפועל בקורס מוגדר בנפרד, בסעיף~\ref{sec:phasea}.

זוהי הרשימה הסגורה של מה שיימסר. תוצר שאינו מופיע כאן אינו חלק מההתחייבות, ותוצר שמופיע כאן ייחשב מסופק רק אם עמד בקריטריון שלו. לגבי גרפים, הסוג נקבע כאן ואינו משתנה בלי בקשת שינוי.

\small

לכל תוצר בטבלה~\ref{tab:deliv} יש קריטריון קבלה, ורק תוצר שעומד בו נחשב כנמסר.

\begin{longtable}{@{}l p{4.6cm} p{2.6cm} p{4.6cm}@{}}
\caption{רשימת התוצרים המחייבת}\label{tab:deliv}\\
\toprule
\textbf{מזהה} & \textbf{תוצר} & \textbf{פורמט} & \textbf{קריטריון קבלה} \\
\midrule
\endfirsthead
\toprule
\textbf{מזהה} & \textbf{תוצר} & \textbf{פורמט} & \textbf{קריטריון קבלה} \\
\midrule
\endhead
\bottomrule
\endfoot
\req{D-01} & פלטפורמת \en{VeriLens} \num{1.0} (שכבות \num{1} עד \num{4}, מסוף, יומן) & שירות ענן וקוד מקור & עמידה ב-\req{T-01} עד \req{T-08} \\
\req{D-02} & תוספים למערכות ניהול תוכן: \en{\mbox{WordPress}} ו-\en{Arc XP} & חבילות התקנה & התקנה והדגמה במערכת פיילוט אחת לפחות \\
\req{D-03} & ממשק \en{API} מתועד, כולל ״ציון שלמות״ & \en{OpenAPI 3.1} & כל נקודות הקצה מתועדות ונבדקו \\
\req{D-04} & מאגר \en{VL-Wild} & \num{20,000} תמונות עם תיוג & תיעוד מקור לכל תמונה; מוחזק אצל צד שלישי \\
\req{D-05} & תקן \en{TRUTH-IMG} גרסה \num{1.0} & מסמך \en{PDF}, עברית ואנגלית & \num{10} כללים, נוהל ביקורת וטופס הסמכה; עבר סבב ציבורי \\
\req{D-06} & לוח בקרה: \textbf{היסטוגרמה} של ציונים, \textbf{עקומת \en{ROC}} חודשית, \textbf{גרף עמודות} של החלטות לפי מחלקה & ממשק \en{Web} & הדגמה (\req{T-08}) \\
\req{D-07} & דוח שקיפות חודשי לדוגמה & \en{PDF} ועמוד באתר & מספרי תמונות, החלטות ותיקונים לחודש פיילוט מלא \\
\req{D-08} & דוח התקדמות & \en{PDF} לפי תבנית הקורס & הוגש בחודש \num{7} \\
\req{D-09} & דוח סיום, כולל תוצאות כל הבדיקות & \en{PDF} לפי תבנית הקורס & מציג עמידה או אי עמידה בכל יעד בטבלה~\ref{tab:goals} \\
\req{D-10} & הדגמה חיה ומצגת & \num{20} דקות & הצגת זרימה מלאה על \num{3} תמונות שיבחר המרצה \\
\end{longtable}
\normalsize

\begin{keybox}[הגדרת ״סיום״]
הפרויקט יושלם כשכל עשרת התוצרים עמדו בקריטריון שלהם. אם יעד מספרי לא הושג, הדוח הסופי יציג את הנתון שהושג בפועל, את הסיבה ואת המשמעות. גם זו תוצאה מחקרית לגיטימית, אבל היא לא תיחשב ״עמידה״.
\end{keybox}

\section{שלב א׳ בקורס: אב טיפוס}\label{sec:phasea}
עד סוף הסמסטר נבנה אב טיפוס קטן שמראה את הרעיון מקצה לקצה, על תמונה אחת: מעלים תמונה, המערכת בודקת תעודת מקור ומריצה גלאי אחד, העורך מחליט, וההחלטה נשמרת ומוצגת לקורא כתווית. אין כאן אימון של מודלים: הגלאי הוא מודל פתוח מוכן, והמטרה היא לבדוק את התהליך ולא לנצח את המתחרים בדיוק. טבלה~\ref{tab:phasea} מגדירה מה יימסר ואיך נבדוק שזה עובד.

\begin{table}[H]\centering
\caption{שלב א׳: מה נבנה בקורס, ואיך נבדוק}\label{tab:phasea}
\small
\begin{tabularx}{\linewidth}{@{}l Y Y@{}}
\toprule
\colhead{מזהה} & \colhead{תוצר} & \colhead{קריטריון קבלה} \\
\midrule
\req{PA-1} & אתר פשוט להעלאת תמונה, שמציג את תוצאות הבדיקה לעורך & מעלים תמונה ומקבלים תוצאה תוך \num{60} שניות לכל היותר \\
\req{PA-2} & בדיקת תעודת \en{C2PA} בעזרת ספרייה פתוחה של \en{C2PA} & מזהה נכון ״חתום / לא חתום / פגום״ בכל \num{20} קובצי הבדיקה \\
\req{PA-3} & גלאי אחד, מודל פתוח מוכן, שנותן ציון לכל תמונה & מדווחים שיעור גילוי ושיעור התרעות שווא על \num{100} תמונות: \num{50} אמיתיות ו-\num{50} שנוצרו בבינה מלאכותית \\
\req{PA-4} & כפתור החלטה לעורך ויומן החלטות & לכל החלטה נשמרים: מי החליט, מתי, הציון והנימוק \\
\req{PA-5} & תווית לקורא, אחת מארבע התוויות (איור~\ref{fig:labels}) & התווית מוצגת ליד התמונה ומקשרת לרשומת ההחלטה \\
\req{PA-6} & מאגר \en{GitHub} עם קוד, \en{README}, פרומפטים ו-\en{commits} של שני השותפים & מישהו אחר מריץ את האב טיפוס לפי ה-\en{README} בלי עזרה \\
\bottomrule
\end{tabularx}
\end{table}

שיעור הגילוי של \req{PA-3} לא חייב לעמוד ביעד של \num{90\%}: זה גלאי אחד ומוכן, ולא המנוע המשוקלל של פרק~\ref{ch:design}. הקריטריון הוא שהמספר נמדד ומדווח בכנות, כולל התמונות שבהן הגלאי טעה.
